\section{Discussion}

\subsection{Key Findings}

\begin{enumerate}
    \item \textbf{The scaling law exists} ($\alpha < 1$), providing a principled alternative to ad-hoc rank selection.
    \item \textbf{The exponent is task-dependent}: $\alpha \approx 0.82$ (single-hop) vs $\alpha \approx 0.30$ (multi-hop).
    \item \textbf{Phase transition at scale}: 12B shows $>$2$\times$ higher rank sensitivity than 1B.
    \item \textbf{Mechanism differs from prediction}: Larger models have \emph{lower} attention entropy.
\end{enumerate}

\subsection{Practical Implications}

Use $r_{\text{opt}} \approx r_{\text{base}} \cdot (N/N_{\text{base}})^\alpha$ as starting point. Calibrate $\alpha$ for your task family.

\textbf{Example} (using $\alpha = 0.82$ for single-hop QA): If $r=16$ works for 7B:
\begin{itemize}
    \item For 70B: $r \approx 16 \cdot (70/7)^{0.82} \approx 106$
    \item For 1B: $r \approx 16 \cdot (1/7)^{0.82} \approx 3$
\end{itemize}

\subsection{Limitations}

\begin{enumerate}
    \item Single architecture (Pythia only)
    \item QA tasks only
    \item Pythia-12B is largest model tested
    \item 3 seeds per configuration
    \item Analysis pipeline validated on synthetic data; full sweeps compute-bound
\end{enumerate}
